\DocumentMetadata{
  pdfversion=2.0,
  pdfstandard=ua-2,
  lang=en-US,
  tagging=on,
  tagging-setup={math/setup=mathml-SE,tabsorder=structure}
}
\documentclass[11pt,letterpaper]{article}
\usepackage[margin=0.68in,includefoot]{geometry}
\usepackage{newcomputermodern}
\usepackage{amsmath,amscd,mathtools}
\providecommand{\square}{\mdlgwhtsquare}
\usepackage[hidelinks]{hyperref}
\hypersetup{
  pdftitle={Analysis First-Year Exam, May 2023, Part 1},
  pdfauthor={University of Florida Department of Mathematics},
  pdfsubject={Graduate mathematics examination},
  pdfdisplaydoctitle=true,
  bookmarksopen=true
}
\setcounter{secnumdepth}{0}
\setlength{\parindent}{0pt}
\setlength{\parskip}{0.28em}
\setlength{\emergencystretch}{3em}
\setlength{\leftmargini}{2.15em}
\setlength{\leftmarginii}{2.65em}
\setlength{\leftmarginiii}{3em}
\renewcommand{\labelenumi}{\textbf{\theenumi.}}
\pagestyle{empty}
\raggedbottom
\allowdisplaybreaks
\newenvironment{examproblems}[1]{%
  \begin{enumerate}%
  \setcounter{enumi}{#1}%
  \setlength{\topsep}{0.35em}%
  \setlength{\partopsep}{0pt}%
  \setlength{\itemsep}{0.48em}%
  \setlength{\parsep}{0.18em}%
}{\end{enumerate}}
\newenvironment{examsubparts}{%
  \begin{enumerate}%
  \setlength{\topsep}{0.18em}%
  \setlength{\partopsep}{0pt}%
  \setlength{\itemsep}{0.16em}%
  \setlength{\parsep}{0.1em}%
}{\end{enumerate}}
\newenvironment{examinstructions}{%
  \par\begingroup\itshape
}{%
  \par\endgroup\vspace{0.18em}
}
\makeatletter
\renewcommand\section{\@startsection{section}{1}{0pt}%
  {-1.25ex plus -.25ex minus -.1ex}%
  {0.55ex plus .1ex}%
  {\normalfont\large\bfseries}}
\renewcommand{\@maketitle}{%
  \newpage\null\vskip -1.2em%
  {\footnotesize\bfseries
    \href{https://www.ufl.edu/}{University of Florida}, \href{https://math.ufl.edu/}{Department of Mathematics}\par}%
  \vskip 0.18em%
  \tagstructbegin{tag=Title}%
    {\LARGE\bfseries\href{https://gma.math.ufl.edu/exams/analysis-first-year/}{Analysis First-Year Exam}, May 2023, Part 1\par}%
  \tagstructend%
  \vskip 0.35em%
  \tagmcbegin{artifact}%
    \hrule height 1pt%
  \tagmcend%
  \vskip 0.55em%
}
\makeatother
\title{Analysis First-Year Exam, May 2023, Part 1}
\author{}
\date{}
\begin{document}
\maketitle
\thispagestyle{empty}
\begin{examinstructions}
Answer FOUR questions in detail. State carefully any results used without proof.
\end{examinstructions}
\begin{examproblems}{0}
\item
Let the continuous function \(f:(0,\infty)\to\mathbb{R}\) satisfy 
\[
\liminf f(n)<0<\limsup f(n).
\]
 Prove that~\(f\) has infinitely many zeros.
\item
Fix any metric space~\(T\) with exactly two points. Prove that the metric space~\(X\) is connected if and only if each continuous function~\(f\) from~\(X\) to~\(T\) is constant.
\item
Let \(B(X)\) be the set of all bounded real-valued functions on the set~\(X\). Prove that the formula 
\[
d(f,g)=\sup\{|f(x)-g(x)|:x\in X\}
\]
 defines a metric~\(d\) on \(B(X)\) that makes \(B(X)\) into a complete metric space.
\item
Let \(A_1\supseteq A_2\supseteq\cdots\) be a decreasing sequence of nonempty subsets of \(\mathbb{R}^N\) (with its standard Euclidean metric).
\begin{examsubparts}
\item[\textnormal{(i)}]
Prove that if each \(A_n\) is compact then \(\bigcap_{n=1}^{\infty}A_n\) is nonempty.
\item[\textnormal{(ii)}]
Show by example that \(\bigcap_{n=1}^{\infty}A_n\) can be empty if each \(A_n\) is only closed.
\end{examsubparts}
\item
Let \(f:(a,b)\to\mathbb{R}\) be differentiable, with bounded derivative. Prove the existence of the real limits 
\[
\lim_{x\to a+}f(x)\quad\text{and}\quad\lim_{x\to b-}f(x).
\]
\end{examproblems}
\end{document}
